\section{One factorial for proof labels and CRT rows}
\label{sec:shared-factorial-carrier}

This section proves the new arithmetic route completely, including its
constructive reverse direction.  Its final Cantor-syntax endpoint is a
quartic in \(349\) nonnegative integer unknowns.  It is the sum of squares
of \(268\) explicitly listed quadratic residuals, or equivalently the
supplied polynomial with \(12\,768\) collected monomials and coefficient
height \(86\,915\,503\).  The square-shell companion has \(351\) natural
unknowns, \(270\) residuals, \(16\,193\) monomials and height
\(15\,317\,516\).

The inherited interface is the literal C222 rule calculus, its existing
ordinary ZFC calibration, and the simultaneous positive-power Pell
representation.  Those results are not reproved or advertised as newly
formalized here.  Everything replacing the old large carrier and CRT
architecture is proved below.  In particular, no incomplete global
encoding is labelled a consistency polynomial.  Neither consistency nor
inconsistency is assumed, and no new end-to-end Lean or object-language
ZFC derivation is claimed.

The structural idea is to reuse \(F=t!\) for two distinct purposes:
the weights \(1+jF\) separate formula-label multiplicities, while the
moduli \(1+F^2/i\) separate the bounded arithmetic rows.  A binomial range
condition applied to a \emph{derived polynomial expression}, rather
than another independently encoded local coordinate, connects the two
roles.  This removes the former enormous carrier power and its fixed
chain of twenty-two squarings.

\subsection{The precise local source interface}

Natural numbers include zero throughout.  Fix a positive horizon \(t\),
outer parameters \(x,y,F\), a row index \(0\le i<t\), and eight local
coordinates
\[
 (a,b,c,d,S,S',u,v).
\]
Their bounds are
\[
 a,b,c,d\le t,\qquad S,S'\le y,\qquad u,v\le x.             \tag{A1}
\]
Let \(c_1,\ldots,c_9\) be the nine literal axiom/schema codes in the
inherited calibrated source, and let \(I,A\) be its implication and
universal constructors.  Both syntax variants use their complete
constructor recipes, not unspecified parameters.  All constructor
expressions are polynomials with nonnegative integer coefficients.
Define
\[
 W_z=1+Fz,\qquad
 G=\bigl((c-a)^2(c-b)^2-(d+1)^2\bigr)^2.
\]
The nonnegative Horn factor is
\begin{align*}
 H={}&(S'-S)^2
       \prod_{\nu=1}^{8}(S'W_{c_\nu}-S)^2\\
    &{}\cdot\bigl((S'W_{c_9}-S)^2+G\bigr)\\
    &{}\cdot(S'W_b-SW_aW_{I(a,b)})^2
       (S'W_{A(a,b)}-SW_b)^2 .
\end{align*}
The local polynomial is
\begin{align}
 L(i,a,b,c,d,S,S',u,v)
   ={}&H+\bigl(x-S-(1+(i+1)y)u\bigr)^2 \notag\\
      &+\bigl(x-S'-(1+(i+2)y)v\bigr)^2.                  \tag{A2}
\end{align}
It is nonnegative at every natural tuple, not only at bounded ones.
The extra expression controlling all possible rule labels is
\[
 J(a,b,c,d)=\sum_{\nu=1}^{9}c_\nu+a+b+I(a,b)+A(a,b)+1.     \tag{A3}
\]
In particular \(J\ge1\).  If \(J\le t\), every label that can occur in
a vanishing Horn branch lies in \([0,t]\).

The finite majorant calculation is performed on these actual source
expressions.  Put
\[
 R=2+t+x+y+F.
\]
Replace every difference square \((r-s)^2\) in the source by
\((U+V)^2\), using positive upper bounds for \(r,s\), and bound each local
input by its corresponding quantity in (A1).  Include both \(L\) and
\(J\), together with \(t+x+y+1\), in the final upper expression.  The
complete gate ledger gives
\[
 0\le L(i,\mathbf w)<R^E,\qquad
 0<J(\mathbf w)<R^E,\qquad E=8\,388\,608.                 \tag{A4}
\]
A separate positive-constructor ledger gives
\[
 J(a,b,c,d)\le(n+1)^\ell
 \quad(0\le a,b,c,d\le n,\ n\ge1),\qquad
 \ell=4\,194\,304.                                      \tag{A5}
\]
These are exact finite certificates.  A constant \(c>0\) is assigned
exponent \(\lceil\log_2 c\rceil\); additions increase the larger
exponent by one; products add exponents.  Since each relevant bounding
base is at least \(2\), every such step is an elementary valid
inequality.  The input bounds and all constructor constants occur in
the ledgers.  Thus (A4)--(A5) do not conceal coefficient factors in an
asymptotic estimate.  Their finite values and the full ledger identities
are independently checked.

\subsection{Reciprocal moduli and the range lemma}

\paragraph{Reciprocal-modulus lemma.}
For \(t\ge1\), set
\[
 F=t!,\qquad T=F^2,\qquad
 d_i=1+\frac{T}{i}\quad(1\le i\le t),\qquad
 B=\binom{T+t}{t}.                                      \tag{A6}
\]
Then every \(d_i\) is an integer greater than one, its prime divisors
exceed \(t\), the \(d_i\) are pairwise coprime, and
\[
 B=\prod_{i=1}^{t}d_i.                                  \tag{A7}
\]

\paragraph{Proof.}
Because \(i\mid F\), one has
\(d_i=1+F(F/i)\equiv1\pmod F\).  A prime at most \(t\)
divides \(F\), and therefore cannot divide \(d_i\).
If a prime divided both \(d_i,d_j\), it would divide
\(id_i-jd_j=i-j\); for distinct \(i,j\), this is impossible
because that prime exceeds \(t\).  Finally,
\[
 \prod_{i=1}^{t}d_i
 =\prod_{i=1}^{t}\frac{T+i}{i}
 =\frac{(T+t)!}{T!\,t!}
 =\binom{T+t}{t}.
\]
The final arithmetic graph materializes only the binomial on the right.
No variable-length product or division by a variable index appears in
the polynomial.  In particular, \(B\ge2\).
\hfill\(\square\)

\paragraph{One-sided prime-power range lemma.}
Let \(p\) be prime, \(e,k\ge1\), and \(a\ge0\).  With the usual convention
\(\binom ak=0\) for \(a<k\),
\[
 p^e\mid\binom ak\quad\Longrightarrow\quad
 a\bmod p^e<k.                                         \tag{A8}
\]

\paragraph{Proof.}
The case \(a<k\) is immediate.  Otherwise the factorial valuation formula is
\[
 v_p\binom ak=
 \sum_{j\ge1}\left(
 \left\lfloor\frac a{p^j}\right\rfloor
 -\left\lfloor\frac {a-k}{p^j}\right\rfloor
 -\left\lfloor\frac k{p^j}\right\rfloor\right).           \tag{A9}
\]
Every summand is \(0\) or \(1\), by the elementary floor-addition identity.
If the interval \([a-k+1,a]\) contained no multiple of \(p^e\),
it would contain no multiple of \(p^j\) for \(j\ge e\).
For those \(j\), its length is less than \(p^j\), and the corresponding
summands in (A9) vanish.  Hence the valuation would be at most \(e-1\),
a contradiction.  Therefore \(p^e\mid a-r\) for some \(0\le r<k\).
The least residue of \(a\) is the least residue of \(r\), which is at
most \(r\), proving (A8).
\hfill\(\square\)

This implication is deliberately one-sided.  For example \(4\bmod4=0<2\),
but \(4\nmid\binom42=6\).  The construction never invokes the false converse.

\subsection{Complete amplified bounded-universal elimination}

Here is the general arithmetic interface used below.  Let \(b_j\ge0\),
let \(t\ge1\), and let
\[
 K\ge\max(1,t,b_1+1,\ldots,b_r+1).
\]
Use (A6) and put \(N=B^K\).  Let \(P(i,\mathbf w)\) be an integer
polynomial, with any fixed external parameters suppressed, such that
\[
 0\le P(i,\mathbf w)<2^K
 \quad\text{for }0\le i<t,\ 0\le w_j\le b_j.              \tag{A10}
\]
Then the following statements are equivalent:
\begin{enumerate}
\item For every \(i<t\), there is a bounded tuple \(\mathbf w\) with
\(P(i,\mathbf w)=0\).
\item There are positive integers \(C,A_1,\ldots,A_r\le N^2\) such that
\begin{align}
 B&\mid C+T+1,                                         \tag{A11}\\
 N&\mid\binom Ct,\qquad
 N\mid\binom{A_j}{b_j+1}\quad(1\le j\le r),              \tag{A12}\\
 N&\mid P(C,\mathbf A).                                \tag{A13}
\end{align}
\end{enumerate}
For soundness the upper bounds \(C,A_j\le N^2\) are not needed.

\paragraph{Soundness.}
Fix \(i\in\{1,\ldots,t\}\), choose a prime \(p\mid d_i\), and write
\(m=v_p(d_i)\ge1\), \(e=Km\).  Pairwise coprimality gives
\(v_p(N)=e\).  Since \(T\equiv-i\pmod p\), (A11) implies
\(C\equiv i-1\pmod p\).  Applying (A8) to the index range gives
\(C\equiv s\pmod{p^e}\) for some \(0\le s<t\).
As \(p>t\), the congruence \(s\equiv i-1\pmod p\) forces \(s=i-1\).
For each local column, (A8) similarly gives
\(A_j\equiv w_j\pmod{p^e}\) with \(0\le w_j\le b_j\).
Polynomial congruence and (A13) now give
\[
 p^e\mid P(i-1,\mathbf w).
\]
But \(p^e\ge2^K\), so (A10) forces this nonnegative value to be zero.
This produces a valid bounded row for each index.

\paragraph{Constructive completeness.}
Select one bounded tuple for each row.  Because the \(d_i\) are pairwise
coprime, CRT supplies representatives in \([1,B^{2K}]\) satisfying
\[
 C\equiv i-1\pmod{d_i^{2K}},\qquad
 A_j\equiv w_{i-1,j}\pmod{d_i^{2K}}.                    \tag{A14}
\]
A zero least residue is replaced by \(B^{2K}\).
Thus the representatives are positive and at most \(N^2\).
Equations (A11) and (A13) follow immediately.

For any requested range lower index \(k\le K\), a representative \(a\)
in (A14) is congruent to some \(w\in[0,k-1]\) modulo \(p^{2Km}\).
If \(a<k\), its binomial is zero.  Otherwise the falling factorial has
the positive factor \(a-w\), divisible by \(p^{2Km}\).
The elementary valuation bound
\[
 v_p(k!)\le k-1\le K-1<Km
\]
shows that division by \(k!\) still leaves valuation at least \(Km\).
This holds for every prime divisor of \(B\), proving all conditions
in (A12).  The strengthening to precision \(2K\) is exactly what
makes completeness valid even when a small prime divides a range
denominator.
\hfill\(\square\)

The free index \(C\) and its range condition are essential.
One must not reuse the old unamplified reciprocal index \(C=B-T-1\).
For \(t=2\), \(F=2\), \(T=4\), \(B=15\), this gives \(C=10\):
it is \(0\) modulo \(5\), but not modulo \(25\).
The frozen construction retains (A11)--(A12).

\subsection{A self-bounding factorial, with both implications}

\paragraph{Factorial lemma.}
Let \(n\ge1\) and let \(f\ge0\) be a proposed factorial output.  Put
\[
 \rho=fn^2,\qquad V=(\rho+1)^n,\qquad D=\binom{\rho+n}{n}.
\]
There exist natural \(u,v\) with
\[
 V+D=Df+u+1,\qquad D=u+v+1                              \tag{A15}
\]
if and only if \(f=n!\).

\paragraph{Proof of soundness.}
The second row gives \(D\ge1\), and subtracting it from the first gives
\(V+v=Df\).  Since \(V\ge1\), one first obtains \(f\ge1\).
The rows say exactly
\[
 f-1<\frac VD\le f,\qquad
 f=\left\lceil\frac VD\right\rceil.                     \tag{A16}
\]
Write \(G=n!\).  The ratio is
\[
 \frac VD
 =G\prod_{i=1}^{n}\frac{\rho+1}{\rho+i}
 =G\prod_{i=1}^{n}
       \left(1-\frac{i-1}{\rho+i}\right).               \tag{A17}
\]
Every factor is positive and at most \(1\).  Moreover
\[
 \sum_{i=1}^{n}\frac{i-1}{\rho+i}
 \le\frac{n(n-1)}{2\rho}
 =\frac{n-1}{2fn}<\frac1{2f},
\]
with the zero sum for \(n=1\) satisfying the same final strict inequality.
The elementary finite-product estimate
\(\prod_i(1-z_i)\ge1-\sum_i z_i\), for \(0\le z_i\le1\),
therefore gives
\[
 G\left(1-\frac1{2f}\right)<\frac VD\le G.               \tag{A18}
\]
Since \(G\) is an integer, (A16) and the upper bound give \(f\le G\).
The lower bound and \(V/D\le f\) give
\[
 G<\frac{2f^2}{2f-1}
   =f+\frac f{2f-1}\le f+1.
\]
Integrality yields \(G\le f\).  Hence \(f=G=n!\).

\paragraph{Constructive completeness.}
Set \(f=G=n!\).  Equation (A18) gives
\(G-\tfrac12<V/D\le G\), so its ceiling is \(G\).
Explicitly define
\[
 u=V+D-DG-1,\qquad v=DG-V.
\]
The ratio inequalities show that these are natural integers, and
direct substitution verifies both equations in (A15).
The proof includes \(n=1\).
\hfill\(\square\)

The output \(f\) occurs in its own radius, but no numerical fixed-point
assumption is made.  It is an existential arithmetic coordinate.
The equations themselves force its unique intended value.

\subsection{The digit graph and a noncircular radix bootstrap}

For clarity, first state the digit interface.  Suppose \(q\ge2\) and
all binomial coefficients \(\binom nj\), \(0\le j\le k\), are less than
\(q\).  If
\[
 (q+1)^n=q^k(qh+c)+r,\qquad
 0\le r<q^k,\quad 0\le c<q,                             \tag{A19}
\]
then \(c=\binom nk\).  This follows by expanding \((q+1)^n\):
the lower coefficients form the carry-free remainder, the coefficient
at \(q^k\) is \(c\), and every higher term is divisible by \(q^{k+1}\).
Conversely the true coefficient supplies the ordinary quotient and
remainder in (A19).

If \(N\mid q\), the fused graph
\[
 (q+1)^n=q^kNh+r,\qquad 0\le r<q^k                       \tag{A20}
\]
is equivalent to \(N\mid\binom nk\), under the same lower-digit bounds.
Indeed the quotient by \(q^k\) is \(\binom nk+qH\).  It is divisible
by \(N\) precisely when the coefficient is.  All reverse witnesses
are nonnegative ordinary quotients and remainders.
The inequalities in (A19)--(A20) are implemented by natural slack rows.

In the final construction \(t\ge16\), and the proposed factorial output
is denoted \(F\).  Define, by ordinary gates and power requests,
\[
 \rho=Ft^2,\quad V=(\rho+1)^t,\quad
 T=F^2,\quad B=\binom{T+t}{t},\quad
 R=2+t+x+y+F,\quad K=2R,\quad N=B^K,
\]
and use one common radix
\[
 q=N(N+V+1)^K.                                         \tag{A21}
\]
At this point the binomial symbols for the proposed \(D,B\) are
\emph{requests}; their intended values must still be justified.
Impose (A15) for \(n=t,f=F\), the raw index row
\[
 C+T+1=Bq_I,                                           \tag{A22}
\]
and the exact/fused digit rows.  The index \(C\) is positive.

The soundness order is as follows.
\begin{enumerate}
\item The raw equation (A22), with \(T=F^2\ge0\), forces \(B>0\).
The base \(\rho+1\) is positive regardless of the proposed \(F\).
The independently proved positive-power graph therefore recovers
\(V\ge1\).
\item The raw ceiling rows give \(D=u+v+1\ge1\) and \(V+v=DF\),
hence \(F\ge1\), without yet interpreting \(D\).
Also \(K>0\).  Recovering \(N=B^K\) gives \(N\ge1\), and recovering
the power in (A21) gives \(q>V^K\ge V^2\).
\item Because \(\rho=Ft^2\ge t\), for \(0\le j\le t\),
\[
 \binom{\rho+t}{j}\le\binom{\rho+t}{t}
 \le(\rho+1)^t=V<q.
\]
The first inequality is monotonicity up to the midpoint; the second
follows termwise from
\(\binom{\rho+t}{t}=\prod_{i=1}^{t}(1+\rho/i)\).
Thus (A19) validates \(D\).  The factorial lemma now gives \(F=t!\).
\item Independently of any intended value of \(B\),
\[
 F^2+t\le(Ft^2+1)^2=(\rho+1)^2.
\]
Hence every coefficient through index \(t\) of its binomial request
is at most \((F^2+t)^t\le V^2<q\).
The second exact digit is valid, so \(B=\binom{F^2+t}{t}\).
Only now use the reciprocal-modulus lemma to conclude \(B\ge2\),
and therefore \(N\ge2^K\).
\end{enumerate}
All remaining power bases are \(q,q+1\) or already recovered positive
quantities.  The simultaneous Pell criterion applies to the finite
power family with its original natural witness conditions.
The recovery does not assume a binomial value before that binomial's
digit has been validated.

For the nine independently chosen representatives (eight columns and
the index), retain the aggregate bound
\[
 C+\sum_{j=1}^{8}A_j+g=9N^2.                            \tag{A23}
\]
Every ordinary range lower index satisfies \(1\le k\le R-1\).
For \(0\le j\le k\), using \(K=2R\),
\[
 \binom nj\le n^k\le9^kN^{2k}
 <2^{4R}N^{2k}\le N^{2k+2}\le N^K<q.                   \tag{A24}
\]
If \(j>n\), the coefficient is zero.  This validates every ordinary
range graph by (A20).  Notice that no interpretation of those ranges
was used to prove their digit bound.

\subsection{Why the amplification coefficient can be small}

The previous amplified route used \(K=8\,388\,608R\).
Here \(K=2R\) suffices because the already required factorial provides
a fixed large floor.  Explicitly,
\[
 16!=20\,922\,789\,888\,000>2^{44}.
\]
For \(E=8\,388\,608\),
\[
 44E<2^{44},\qquad 2E\le2^{44}.                         \tag{A25}
\]
At \(r_0=2^{44}\), one has \(r_0^E=2^{44E}<2^{r_0}\).
For any integer \(r\ge r_0\),
\[
 \left(1+\frac1r\right)^E
 \le\sum_{j=0}^{E}\left(\frac Er\right)^j<2,
\]
since \(E/r\le1/2\) and the geometric sum is finite.
If \(r^E<2^r\), multiplication by this inequality gives
\((r+1)^E<2^{r+1}\).  Integer induction proves
\[
 r^E<2^r\quad(r\ge2^{44}).
\]
After the factorial bootstrap, \(R\ge F=t!\ge16!\), so
\[
 L(i,\mathbf w),J(\mathbf w)<R^E<2^R<2^K.               \tag{A26}
\]
This is exactly the unbounded evaluation bound needed for CRT decoding.
It is used only \emph{after} the factorial has been recovered, not as a
circular premise for the earlier digit interpretation.

The preserved hybrid comparison admits another premise.
If an existing carrier gives
\(R\ge2^\lambda\), \(\lambda\ge8\), \(E\le2\lambda\), then
\(2^\lambda>2\lambda^2\ge\lambda E\) and \(2^\lambda\ge2E\).
The same base case and induction prove \(R^E<2^R\).
The final shared-factorial endpoint does not retain that old carrier;
it uses the explicit factorial floor (A25).

\subsection{Complete factorial-weight Horn-flow equivalence}

\paragraph{Weight separation.}
For \(1\le j\le t\), put \(W_j=1+jF\), \(F=t!\).
These weights are pairwise coprime.  If a prime \(p\) divided
\(W_j,W_k\), it would divide \((j-k)F\) but not \(F\).
Thus \(p\mid j-k\), so \(p\le t\) and \(p\mid F\), a contradiction.
Every positive-label weight is greater than one and has a prime divisor.
The truth weight is \(W_0=1\).

\paragraph{Soundness of the encoded path.}
Suppose each row \(i<t\) has a bounded zero of (A2), with
\(J(a,b,c,d)\le t\).  The beta equalities give
\[
 x=S_i+(1+(i+1)y)u_i,\quad
 x=S'_i+(1+(i+2)y)v_i.
\]
As \(S_i,S'_i\le y\), uniqueness of remainders modulo the common
overlap modulus shows \(S'_i=S_{i+1}\) for \(i<t-1\).
The outer endpoint rows are
\[
 x=F+1+(y+1)u_0,\qquad y=F+1+r_0,\qquad
 x=1+(1+(t+1)y)u_t.                                    \tag{A27}
\]
They force \(S_0=F+1=W_1\) and \(S'_{t-1}=1\), because these
also lie in the prescribed remainder intervals.

Every summand in (A2) is nonnegative.  Thus its Horn product vanishes,
and at least one actual rule branch vanishes.  Choose such a branch
and denote its premises by \(p_i,q_i\) and conclusion by \(v_i\),
using truth label \(0\) for unused premises.  The branch equation is
\[
 S'_i W_{v_i}=S_iW_{p_i}W_{q_i}.                        \tag{A28}
\]
For an axiom both premises are \(0\); for generalization one is \(0\).
The truth branch has all three labels \(0\).
The bundled schema factor is a sum of two squares, so its vanishing
also enforces its actual side condition.  A purely numerical match
with an invalid schema instance is not accepted.

All labels in (A28) are in \([0,t]\), by the definition of \(J\).
All weights are positive.  Starting from \(S_0=W_1>0\), equation (A28)
inductively forces every \(S'_i\), hence every \(S_i\), positive.
Multiplying the row equalities and cancelling the overlapping positive
states is therefore legitimate.  The endpoints yield
\[
 \prod_i W_{v_i}=W_1\prod_iW_{p_i}W_{q_i}.               \tag{A29}
\]
For any positive label \(j\le t\), choose a prime \(p_j\mid W_j\).
Pairwise coprimality annihilates the valuations of every other weight.
Taking the \(p_j\)-valuation in (A29) and cancelling the positive integer
\(v_{p_j}(W_j)\) gives
\[
 \#\{i:v_i=j\}
 -\#\{i:p_i=j\}-\#\{i:q_i=j\}=\delta_{j1}.              \tag{A30}
\]
The weights need not be prime or square-free.

Let \(D_0\) be the closure of truth under the finitely many chosen valid
rules.  If falsity \(1\) were outside \(D_0\), each rule with conclusion
outside \(D_0\) would have at least one premise outside \(D_0\).
Hence the total number of outside conclusions would be at most the
number of outside premises.  Summing (A30) over the positive labels
outside \(D_0\), however, gives the difference equal to \(1\).
This contradiction proves \(1\in D_0\), hence a finite actual
refutation.  This argument admits cyclic, repeated and unordered rule
occurrences.  No fictitious balance equation for label \(0\) is used:
truth is already in \(D_0\).

\paragraph{Constructive completeness.}
Start from a finite C222 derivation of falsity and expand its proof DAG
to a finite proof tree.  Collect its finitely many parameter tuples and
compute the full value of \(J\) for each tuple, including the codes of
branches not used by that tuple.  Also include the fixed truth-padding
tuple \(a=b=c=d=0\).  First choose \(t\) at least \(16\), at least the
number of tree steps, at least every parameter, and at least every one
of those finitely many label sums.  This choice does not involve \(F\)
or any accumulator.  Only then define \(F=t!\) and the weights.

Start with the pending multiset containing falsity.  Expand a proof-tree
conclusion \(v\) backwards to its premises \(p,q\).  The product of
weights of the pending multiset before the expansion is \(S\);
the product after it is \(S'\).  Both are positive integers and satisfy
(A28).  Truth leaves have unit weight.  Once the proof tree is
exhausted the product is \(1\).  Insert truth steps to obtain exactly
\(t\) rows; they leave the current product unchanged and their label
sums were included in the horizon choice.

Now choose \(y\) divisible by \(t!\), at least every accumulator,
and at least \(F+1\).  The \(t+1\) moduli
\(1+(i+1)y\), \(0\le i\le t\), are pairwise coprime:
a common prime would divide an index difference at most \(t\) while
remaining coprime to \(y\), impossible because \(t!\mid y\).
CRT gives a nonnegative \(x\) with all prescribed state residues.
Every residue lies in \([0,y]\), below its modulus, so the corresponding
Euclidean quotients are natural and at most \(x\).
The outer slacks and quotients in (A27) are natural as well.
This constructs the bounded eight-column source path.  In particular,
the horizon is not required to dominate the accumulator sizes,
which are chosen and encoded only afterwards.
\hfill\(\square\)

\subsection{Derived polynomial ranges and their independent carry bound}

\paragraph{General derived-range lemma.}
Use the amplified CRT construction above and now require \(K\ge t+1\).
Let \(J\) be an integer polynomial nonnegative on natural inputs, with
\(0\le J(\mathbf w)<2^K\) for every bounded local tuple.
Adding
\[
 N\mid\binom{J(\mathbf A)}{t+1}                         \tag{A31}
\]
forces \(J(\mathbf w)\le t\) in every decoded row.

Indeed, for the relevant prime power \(p^e\mid N\), (A8) gives
\(J(\mathbf A)\equiv r\pmod{p^e}\) for some \(0\le r\le t\).
Polynomial congruence gives \(J(\mathbf w)\equiv r\pmod{p^e}\).
Because \(J(\mathbf w)<2^K\le p^e\) and \(t<K\le p^e\),
these are both least nonnegative residues.  Thus \(J(\mathbf w)=r\le t\).

Conversely, suppose the selected local tuples have \(J(\mathbf w)\le t\).
The precision-\(2K\) representatives give
\(J(\mathbf A)\equiv J(\mathbf w)\pmod{p^{2Km}}\).
If \(J(\mathbf A)<t+1\), its binomial is zero.  Otherwise its falling
factorial contains a factor divisible by \(p^{2Km}\).
Division by \((t+1)!\) loses at most \(t<K\le Km\) in valuation,
so (A31) follows for every prime of \(B\).
This is the complete reverse direction and does not require any
independent CRT coordinate for \(J\).
\hfill\(\square\)

For the specific label expression (A3), the needed local bound is
(A26).  Thus (A31) replaces the extra witness \(w\) and residual
\((t-w-J)^2\) of the nine-column predecessor.  The local source
(A2) itself is unchanged.

There is an important distinct implementation obligation:
the digit graph for (A31) must be carry-free even though its upper
argument is a derived polynomial.  The aggregate (A23) must
\emph{not} include \(J(\mathbf A)\), which may be much larger than
\(N^2\).  The ordinary CRT construction cannot promise such an
aggregate bound for a derived value.

Let \(n\) be the largest of the four encoded constructor parameters.
By (A23), \(n\le9N^2\); because \(N\ge2^K\) and \(t\ge16\),
\(N\ge10\).  Hence
\[
 n+1\le9N^2+1\le N^3.
\]
The independently certified positive-DAG bound (A5) gives
\[
 J(\mathbf A)\le(n+1)^\ell\le N^{3\ell}.                \tag{A32}
\]
The finite numerical base check is
\[
 2\cdot16!\ge3\ell\cdot17.
\]
It persists for every \(t\ge16\): if
\(2t!\ge3\ell(t+1)\), then
\[
 2(t+1)!\ge3\ell(t+1)^2\ge3\ell(t+2).
\]
Consequently
\[
 3\ell(t+1)\le2F<2R=K.                                \tag{A33}
\]
For every lower digit \(0\le j\le t+1\), (A32)--(A33) yield
\[
 \binom{J(\mathbf A)}j
 \le J(\mathbf A)^{t+1}
 \le N^{3\ell(t+1)}
 \le N^K<q.                                           \tag{A34}
\]
Again a coefficient above the upper index is zero.
This proves that the additional fused digit graph recognizes exactly
(A31).  Its bound is independent of the representative aggregate and
explicitly accounts for the entire constructor degree and coefficients.
There are ten range requests in total, but only nine independent
representatives in (A23).

\subsection{Assembly into one complete polynomial}

The finite arithmetic system now consists of the outer and endpoint
rows, the single self-bounding factorial, the reciprocal product
binomial, the amplification and common radix, nine ordinary ranges,
one derived-label range, the index congruence and the local evaluation
divisibility.  Every power request is replaced by the inherited explicit
simultaneous Pell graph.  Every ordinary gate is quadraticized while
sharing identical intermediate expressions.  If the resulting quadratic
residuals are \(r_1,\ldots,r_s\), define
\[
 Q(\mathbf z)=\sum_{\nu=1}^{s}r_\nu(\mathbf z)^2.         \tag{A35}
\]
Its complete literal rows and expanded coefficients are supplied.
For natural inputs, \(Q=0\) is equivalent to every row being zero.

\paragraph{Soundness of the assembled equation.}
From a natural zero, recover the true powers and interpret the exact
factorial and reciprocal-product digits in the order proved above.
Then use the factorial floor to establish (A26), and the separate
ordinary and derived carry estimates to interpret all range graphs.
Amplified CRT decoding yields a bounded zero of (A2) at every row,
and the derived-range lemma gives \(J\le t\) there.
The factorial-weight Horn-flow theorem produces a finite C222
refutation.  Applying the inherited source calibration gives the stated
ZFC refutability predicate.

\paragraph{Constructive completeness of the assembled equation.}
Start with a finite calibrated refutation.  Choose the padded horizon
and construct the shared-factorial path before choosing its beta
encoding, as proved above.  Set \(F=t!\), \(T=F^2\),
\(B=\binom{T+t}{t}\), \(K=2R\), and \(N=B^K\).
Choose the eight columns and the index by precision-\(2K\) CRT.
Their positive representatives are at most \(N^2\), so (A23) has a
natural slack.  The local polynomial remains nonnegative on all
natural inputs; its divisibility quotient is therefore natural.
The derived-label lemma supplies (A31).
Define the true \(\rho,V,q\); the factorial lemma supplies its two
ceiling witnesses, and the carry-free estimates supply the actual
digit quotients and remainders.  Finally the existing simultaneous
Pell completeness theorem supplies one common frame and all witnesses
for this finite family of true positive powers.  Thus the full
quadratic system, and hence (A35), has a natural zero.

The final nonnegative contraction deletes coordinates only with exact
integral reverse substitutions certified to be nonnegative on the
remaining natural variables.  For signed-coefficient restorations,
the certificate explicitly expands a positive rational sum of squares
of affine forms plus a coefficientwise nonnegative remainder.
Every identity and substitution is independently replayed, so the
contraction preserves both zero-existence directions, not merely the
forward map.  It leaves \(349\) variables and \(268\) quadratic rows
in the Cantor endpoint.  Independent expansion of their squares gives
the stated exact quartic and coefficient ledger.

\paragraph{Measured tradeoffs and finite evidence.}
The new endpoint uses one factorial, two exact binomials, ten range
graphs, nineteen positive-power requests, sixteen independently tested
Pell indices and one derived index.  Normal and Python-optimized
regeneration reproduce all twelve final endpoint files byte for byte.
The semantic regressions include \(361\,200\) range cases, \(32\)
constructive CRT cases, ten derived-range constructions, \(600\)
factorial candidates and \(47\,904\) small Horn-rule multisets.
Negative controls expose the false converse to (A8), invalid canonical
index lifting, a bad CRT row, missing label bounds, and derived values
that exceed the representative aggregate.  These finite tests support
the implementation; the quantified proofs above establish the
unbounded implications.

The \(350\)-variable local-slack predecessor has fewer monomials than
the \(349\)-variable derived-range endpoint.  The earlier \(363\)-variable
reciprocal construction has fewer still.  These are explicit
variable/support/height tradeoffs; no global optimum or worldwide
priority is asserted.  All inherited code, calculus, Pell and
contraction dependencies retain their original provenance and notices.
